\maketitle
\begin{abstract}
The Bohigas--Giannoni--Schmit (BGS) conjecture relates the microscopic
spectra of appropriately desymmetrized quantum chaotic systems to
random-matrix universality. The original chaotic-billiard evidence concerned
the orthogonal symmetry class \cite{BGS1984}; the specified unitary-class
benchmark in this programme is the $\mathrm{Sine}_2$ process. This
\emph{continuing research monograph} collects distinct, rigorously
qualified project results and their proof boundaries. It does not establish
BGS for the fixed deterministic map, let alone for all chaotic systems.

The arithmetic programme develops finite-field pair-key and Cayley-energy
reductions, cyclotomic rigidity, exact quantum-map multiplicities,
arithmetic-subsequence return laws, and smooth/Gevrey cohomological and
resonance results. Some statements concern compressed observables or
vanishing classical kicks; others concern fixed-strength but specially
conjugate quantum systems. Related chapters address locality, harmonic
analysis, spectral obstructions and quantum-information estimation.
Their hypotheses, symmetry classes, time scales and spectral objects
remain separate, rather than being combined into a universality claim.

The central deterministic problem uses the odd-dimensional family
$U_N=\exp(i\tau N A_N)M_N$ with $\tau=10^{-80}$ \emph{fixed as $N\to\infty$}.
The manuscript establishes exact trace/resolvent identities, conditional
correlation bridges, quantitative intensity and selected interference
estimates, and auxiliary noisy comparisons. An internally reconstructed
continuation gives phase-memory bounds, an absorbed scalar real-flow
comparison, and the sharp rational matrix-inner second-jet constant
$4/3$. The actual all-order, signed Heisenberg-scale trace comparison,
including the pair plateau correction and deterministic removal of
auxiliary noise, remains \textbf{open}. The reported proofs have not
received a blanket independent external referee or novelty certification.
\end{abstract}
\vspace{1ex}
\noindent\textbf{Keywords.} BGS conjecture; quantum cat maps; Sine$_2$;
signed joint traces; Weyl quantization; arithmetic return laws;
cohomological resonance; exact certificates.\\
\textbf{Mathematics Subject Classification.} 81Q50, 37D20, 11T23, 11R18.
{\small\tableofcontents}
\clearpage

\section{Scope, normalizations, and the dependency structure}
\label{sec:scope}
This is one consolidated research manuscript, not a concatenation of claims
of equal evidentiary strength. Historical versions remain archived. Repeated
proofs of the same mechanism have been combined or cross-referenced; distinct
asymptotic regimes have not been identified. In particular, a global Gaussian
law of an averaged observable and a sine-kernel law for local propagator
levels are different assertions. The former supplies neither the latter
nor Heisenberg-scale joint trace control.

\subsection{The deterministic microscopic target}
\label{subsec:unresolved-deterministic-target}
In this manuscript the positive unitary-class benchmark is an
explicitly prescribed quantum torus family, rather than the original
GOE-class Sinai-billiard calculation \cite{BGS1984}. The maintained
frozen propagator is
\[
 S=\begin{pmatrix}3&4\\2&3\end{pmatrix},\qquad
 U_N=\exp(i\tau N A_N)M_N,\qquad
 \tau=10^{-80},\quad N\ge3\text{ odd},
 \label{eq:frozen-target-introduction}
\]
where $M_N$ is the specified torus quantization and $A_N$ is the
specified Weyl-quantized kick. The coupling $\tau$ is \emph{independent}
of $N$; letting it tend to zero, averaging independent diagonal phases,
or adding stochastic noise studies a different family.

A necessary CUE pair-reference target in the maintained normalization is
\[
 L_N(s)=N^{-2}\sum_{m\ge1}e^{-2sm/N}
       |\operatorname{Tr}U_N^m|^2
 \longrightarrow \frac{1-e^{-2s}}{4s^2},
 \qquad s>0.
 \label{eq:frozen-pair-abel-introduction}
\]
It has \emph{not} been proved for this deterministic map. Even proving
this pair identity would not suffice for convergence of every fixed
factorial order: the missing assertion is the fully
\emph{signed, symmetry-resolved Heisenberg-scale joint-trace
comparison}, with the modulus taken outside each complete Fourier
sum, together with the necessary local-process tightness conditions.
For the ordinary unitary bulk, characterization through complete
loop-equation hierarchies is another possible conditional route
\cite{BH2026}; the required model-specific approximate equations
and source bounds have not been established here. No conclusion about
the unchanged $U_N$ is imported from that external characterization.

\subsection{The three spectral objects}
We distinguish a quantum propagator $U_N$, a compressed observable
$P_jB_NP_j|_{\operatorname{ran}P_j}$, and the matrix elements of $B_N$ in an
eigenbasis. Universality for one of these is not universality for another.
An empirical law of $U_N^K$ folds its phases K times; it is not the original
nearest-neighbour law. A two-scale measure of one eigenvalue is not a
two-eigenvalue correlation. These distinctions are used throughout.

In the finite-field sections, p is an odd prime, H is a multiplicative
subgroup and K is its order. In the Pell sections, N is the Hilbert-space
parameter and can be composite, while K is an arithmetic return time with
$N\asymp(3+2\sqrt2)^{K/2}$. No unexplained identification of p and N is made.
The notation D denotes the dimension of the active parity sector except
where an operator is explicitly introduced under another local convention.
Normalized Hilbert--Schmidt norms are $\|X\|_{2,D}^2=D^{-1}\operatorname{Tr}X^*X$;
operator norms carry the subscript $\mathrm{op}$. Replacing one by the other
can change a valid statement into a false one.

\subsection{Weyl, classical, and perturbation conventions}
For the fixed Fourier-covector matrix
$T=\left(\begin{smallmatrix}3&2\\4&3\end{smallmatrix}\right)$, the
classical pullback matrix is $S=T^{\mathsf T}$. Define
\[
 W_N(a,b)\psi(x)=e_N(a(x-b/2))\psi(x-b),\qquad
 Q_N(a,b)=(-1)^{ab}W_N(a,b),\qquad\hbar_N=(2\pi N)^{-1}.
\]
The half in W is modular; a,b in Q are integer Fourier frequencies before
reduction. Thus $Q_N(v)Q_N(w)=e^{\pi i\omega(v,w)/N}Q_N(v+w)$.
For general hyperbolic matrices an explicitly stated invariant characteristic
$Q_N^*(a,b)=(-1)^{ab+a+b}W_N(a,b)$ is used instead.
These characteristics are not silently interchanged.

A quantum kick $e^{i\varepsilon\operatorname{Op}_N(f)}$ has classical action
strength $\kappa=\hbar_N\varepsilon$. The fixed-strength convention in the
new return-word section is $e^{i\kappa\operatorname{Op}_N(f)/\hbar_N}$.
The older locality section writes $e^{i\kappa_0N\operatorname{Op}_N(f)}$;
there $\kappa=\kappa_0/(2\pi)$. Its formulas retain the explicitly stated
$\kappa_0$ convention. This factor does not disappear on consolidation.
We use $\{f,h\}=f_qh_p-f_ph_q$ and $X_f=(f_p,-f_q)$; sources using
$(-f_p,f_q)$ have their flow-time sign stated alongside the map.

\subsection{What is proved, imported, checked, or unresolved}
The internal dependency chains are as follows. Finite Weyl covariance gives
fixed-point trace identities and the exact cyclic multiplicities. Quadratic
norm lifting, bounded-cluster counting and summable Pell phases give the
observable moment laws. A finite Duhamel estimate then transfers selected
laws to recurrence times of the actual propagator. The much stronger
fixed-classical-strength results use support locality or noncommuting
return words; they do not follow from the moment laws.

For the second resonance, the orbit expression of Green--Kubo variance gives
its exact zero criterion. A local invariant differential operator gives the
Hessian obstruction. The finite-box theorem additionally depends on an exact
proof object, reconstructed and replayed by two implementations; it is not
inferred from random sampling. The return-word example proves that a nonzero
bracket can still lie in that zero set on the smooth infinite-frequency class.
No implication from analytic nonzero bracket to positive variance is inserted.

The finite-field reductions form a parallel chain: exact pair keys, Cayley
energy, fixed-degree row polynomials, scalar supports, and conditional
arithmetic estimates. They are not a substitute for spectral trace estimates.
The subresultant, cyclotomic-unit and subgroup-curve inputs cited in those
sections retain their named external hypotheses. Their historical derivations
are preserved in this consolidated record; the present audit has not supplied
new proofs of all external results. Exact-source checks and unretrieved leaves
are listed separately in the source ledger, so no blanket referee verdict is
attached to the whole historical corpus.

The scalar Fejer--Riesz theorem is used only after strict positivity of the
proposed Laurent polynomial has been proved. The classical orbit/coboundary
criterion is consistent with Cohen--Conze~\cite{CohenConze}; it is not
imported as a quantum central limit theorem. Residual perturbations and their
second-order commutator have direct precedent in Prosen--Znidaric~\cite{PZ2005}.
Nonlinear cat-map pseudo-symmetries also predate this project~\cite{KM2000}.
The finite-return identity is group-valued coboundary algebra; no priority
claim for that algebra is made. All new conclusions below are derived directly
rather than assigned to a source merely from its title or abstract.

\subsection{A precise universality obligation}
For a D-dimensional spectrum $e^{2\pi ix_a}$, the global-density-unfolded
factorial pair statistic is
\[
 R_D(\phi)=D^{-1}\sum_{a\ne b}\sum_{m\in\mathbb Z}
                   \phi(D(x_a-x_b+m)).
\]
Poisson summation yields
\[
 R_D(\phi)=D^{-2}\sum_{k\in\mathbb Z}\widehat\phi(k/D)
                  (|\operatorname{Tr}U_D^k|^2-D).
\]
Thus a positive local-universality result requires the appropriate weighted
joint trace comparison on the Heisenberg scale, with symmetry reduction and
unfolding fixed beforehand. Positivity can sometimes exclude a proposed limit
using only a few shorter-time traces; it cannot prove the full comparison.
The two principal unresolved obligations remain a sharp arithmetic closure
in the split finite-field branch and positive spectral universality for a
natural fixed map not retaining the arithmetic model by smooth conjugacy.
Neither is declared solved by the present consolidation.

\subsection{Reading map for the continuing programme}
\label{subsec:reader-map}
The order of chapters records related research developments, not a single
chain of implications to microscopic universality. For orientation the
mathematical layers should be read with the following distinctions:
\begin{itemize}[leftmargin=*,itemsep=2pt]
\item \textbf{Finite-field and Cayley arithmetic:} exact combinatorial
reductions and conditional or range-restricted energy bounds; the sharp
centered pointwise arithmetic target remains open.
\item \textbf{Return and Pell/Floquet theory:} asymptotic laws on explicitly
declared arithmetic subsequences, regularity scales and observables; these
cannot be silently reinterpreted as microscopic level statistics.
\item \textbf{Resonance and cohomology:} sharp coercivity, rigidity and
counterexamples in distinct smooth, analytic, linear-transport and nonlinear
regimes; the all-members analytic extension remains unresolved.
\item \textbf{Frozen-map microscopic programme:} symmetry, exact trace
identities, selected interference sectors and conditional local-law
criteria; the signed Heisenberg-scale comparison at every fixed factorial
order, including its pair-scale plateau correction, is still missing.
\item \textbf{Separate information-theory and noisy-law branches:} theorems
with their own dynamics and source/measurement assumptions; these are not
results about the unchanged deterministic fixed-strength spectrum.
\end{itemize}
The authoritative status and dependencies are the maintained claim ledger
and mathematical source files in the companion repository. A section's
historical title, independent finite replay or successful typesetting run
does not revise those statuses. This long-form manuscript is a living
research exposition; focused possible articles have separate claim
manifests and publication-review requirements.
